% ============================================================
% 2_B_template.tex --- 管理员模板管理功能需求
% 编写人：B  日期：2026-07-15
% 职责模块：学院管理 / 模板管理 / 论文CRUD / 章节存取
% ============================================================

\section{管理员模板管理功能需求}

模板是论文自动生成系统的核心基础，定义了论文的封面样式、字体字号、页边距、页眉页脚、目录格式等全部排版规范参数。管理员通过模板管理功能，维护不同学位类型（本科毕业论文、课程论文、项目论文）、不同学院维度的论文模板，确保所有学生生成的论文文档严格符合学校教务处的排版要求。

\subsection{模板管理功能模块划分}

模板管理功能划分为以下五个子模块：

\begin{enumerate}
    \item \textbf{模板基础管理}：模板元数据的创建、编辑、删除和列表查询；
    \item \textbf{模板版本管理}：模板配置变更时自动生成新版本，保留完整历史记录，支持版本切换与回滚；
    \item \textbf{模板排版参数配置}：封面元素、章节结构、字体格式、页边距、页眉页脚等全部排版参数的在线可视化配置；
    \item \textbf{模板预览与校验}：管理员配置模板后实时预览生成效果，系统自动校验配置合法性；
    \item \textbf{学院管理}：维护学校各学院的基础信息（名称、代码），模板按学院维度进行关联绑定。
\end{enumerate}

\subsection{模板基础管理}

\reqitem{FR-B01}{管理员可创建论文模板，填写模板名称（如\textquotedblleft 计算机科学与技术学院本科毕业论文模板\textquotedblright），选择模板类型（GRADUATION 毕业设计 / COURSE 课程论文 / PROJECT 项目报告），选择适用学院（支持多选，也可设为\textquotedblleft 全部学院\textquotedblright 通用模板）。模板创建后自动生成初始版本（V1.0），初始配置为空，需管理员进一步编辑排版参数。}
\reqitem{FR-B02}{管理员可编辑已有模板的元数据（名称、适用学位类型、适用学院范围）。编辑元数据不触发版本号变更；仅当修改排版配置参数时系统自动递增版本号。}
\reqitem{FR-B03}{管理员可删除不再使用的模板。删除操作前，系统自动检测该模板的历史版本是否被现有论文引用：若存在关联的进行中论文（status 不为 COMPLETED），提示\textquotedblleft 该模板下有 N 篇进行中论文，无法删除\textquotedblright 并阻止删除；若无关联论文或仅有已完成论文，允许删除，同时删除所有关联的 TemplateVersion 记录。}
\reqitem{FR-B04}{管理员进入模板列表页面，支持按模板类型（GRADUATION / COURSE / PROJECT）和学院进行筛选。列表展示字段包括：模板名称、模板类型、适用学院、当前生效版本号、创建时间、最后修改时间。点击模板行进入版本管理详情页。}

\subsection{模板版本管理}

\reqitem{FR-B05}{每次管理员修改模板的排版配置（coverFields / chapterStructure / formatConfig 中任一 JSON 字段）并保存时，系统自动执行版本管理流程：将当前生效版本 is\_current 置为 false；解析当前版本号递增小版本（V1.0\textrightarrow V1.1，V1.9\textrightarrow V2.0）；创建新 TemplateVersion 记录并将 is\_current 置为 true。修改模板元数据（名称、类型）不触发版本递增。}
\reqitem{FR-B06}{管理员可在版本历史页面查看某一模板的全部历史版本列表，展示版本号、创建时间、是否为当前生效版本。支持将任一历史版本设为当前生效版本（回滚操作）：将目标版本的 is\_current 置为 true，原当前版本置为 false。回滚操作不影响已使用该模板的历史论文（论文绑定的 template\_version\_id 保持不变）。}
\reqitem{FR-B07}{管理员选择两个历史版本进行差异对比，系统逐字段比对 coverFields、chapterStructure、formatConfig 三个 JSON 字段的内容差异，前端以对比视图展示变更内容：新增部分以绿色高亮标记，删除部分以红色高亮标记，未变更部分灰色显示。}

\subsection{模板排版参数配置}

\reqitem{FR-B08}{管理员进入模板版本编辑页面，可配置页面参数：纸张大小（默认 A4）、上/下/左/右边距（单位 cm）、行间距倍数（1.5 倍 / 2.0 倍）、段落间距（单位 em）。修改后右侧预览区域实时刷新渲染效果。}
\reqitem{FR-B09}{管理员可配置字体参数：正文字体与字号（默认宋体 12pt）、一级标题字体与字号（默认黑体 16pt）、二级标题字体与字号（默认黑体 14pt）、三级标题字体与字号（默认黑体 12pt）、页眉页脚字体与字号。字体列表从系统可用字体中加载。}
\reqitem{FR-B10}{管理员通过可视化编辑器配置封面元素：校徽图片（上传+预览+位置拖拽）、论文标题（字体/字号/水平位置/垂直位置）、作者姓名格式、指导教师信息格式、学院名称、提交日期格式（支持 yyyy年MM月dd日 等自定义格式）。每个元素独立可拖拽调整位置（x/y 百分比坐标）。}
\reqitem{FR-B11}{管理员配置页眉页脚样式：奇数页页眉内容（可选\textquotedblleft 章标题\textquotedblright / \textquotedblleft 论文题目\textquotedblright / 自定义文本）、偶数页页眉内容、页脚页码格式（如\textquotedblleft 第\{page\}页\textquotedblright 或\textquotedblleft -\{page\}-\textquotedblright）、页码位置（居中/居右）。}
\reqitem{FR-B12}{管理员通过章节树编辑器配置模板的章节结构（chapterStructure）：以树形结构展示论文章节骨架，支持添加/删除章节节点、拖拽调整章节层级和顺序、编辑章节标题。一章节点下可包含多个子节节点，子节节点可继续嵌套。模板内置默认章节结构（摘要、目录、引言、正文各章、结论、致谢、参考文献），管理员可在此基础上自定义调整。}

\subsection{模板预览与校验}

\reqitem{FR-B13}{管理员配置模板排版参数后，可点击\textquotedblleft 在线预览\textquotedblright 按钮。系统使用预设的示例论文内容（包含中英文摘要、目录、多级标题正文、图片、表格、参考文献等代表性元素），按照当前模板配置渲染生成完整的预览 PDF。预览 PDF 在浏览器新标签页中展示，支持缩放和翻页。}
\reqitem{FR-B14}{保存模板配置前，系统自动执行合法性校验，校验项包括：封面必填元素（标题、作者、日期）是否已配置；章节结构是否至少包含\textquotedblleft 正文\textquotedblright 类型的章节节点；字体名称是否在系统支持范围内；页边距设置是否在打印机可打印区域内（通常上下左右均不低于 1.5cm）；JSON 字段格式是否正确。校验不通过时，以红色错误提示框列出所有不合规项及其修复建议，阻止保存。}

\subsection{学院管理}

\reqitem{FR-B15}{管理员可新增学院信息，填写学院名称（如\textquotedblleft 计算机科学与技术学院\textquotedblright）和学院代码（如\textquotedblleft CS\textquotedblright，全局唯一，用于系统内部标识和API路由）。学院代码一经创建后不可修改，但学院名称可编辑。}
\reqitem{FR-B16}{管理员可编辑学院名称，或删除学院。删除操作前，系统自动检测该学院是否存在关联的模板或论文记录：若存在关联模板，提示\textquotedblleft 该学院下存在 N 个关联模板，请先处理关联模板后再删除\textquotedblright；若存在进行中论文，提示\textquotedblleft 该学院下存在进行中的论文，无法删除\textquotedblright。无关联数据时允许删除。}
\reqitem{FR-B17}{管理员在创建模板选择适用学院时，从已维护的学院列表中通过多选下拉框选择。支持\textquotedblleft 全部学院\textquotedblright 快捷选项，选中后该模板适用于所有学院（college\_id 设为 NULL）。}

